% 根据课题选择试验、模型、工程调查等方法，不必保留不适用的小节。
\chapter{研究对象与方法}
\section{研究对象与边界条件}
本节说明工程对象、采样条件、装置规模与运行工况，明确研究边界及试验条件与实际工程条件的关系。
\section{实验设计与分析方法}
\subsection{实验方案与对照设置}
本节说明试验分组、对照条件、平行样数量、运行周期及数据处理方法，并区分独立重复与重复测定。
\subsection{质量保证与质量控制}
本节说明检测依据、仪器设备、校准方法及质量控制措施，必要时列明检出限、定量限和数据有效性判据。
\section{指标计算与单位表达}
物理量采用斜体，单位及描述性下标采用正体。公式按章连续编号，编号置于右侧。以浓度为依据计算去除率时，可表示为
\begin{equation}
  \eta=\frac{C_{\mathrm{in}}-C_{\mathrm{out}}}{C_{\mathrm{in}}}\times100\%
  \label{eq:removal}
\end{equation}
式中，$C_{\mathrm{in}}$与$C_{\mathrm{out}}$分别为进、出水目标物质量浓度，单位为\unit{\milli\gram\per\litre}。使用该指标时，应说明采样条件及浓度比较的适用前提。
% 流量变化明显时需考虑质量通量，不能把浓度去除率直接当作质量去除率。
% 常用写法：\qty{20}{\degreeCelsius}、\unit{\kilo\watt\hour\per\metre\cubed}。
% 化学式：\ce{NH4+}、\ce{NO3-}；pH 正体：\(\mathrm{pH}\)。
\subsection{统计分析与不确定性}
本节说明统计方法、样本量、误差表示方式、模型假设及不确定性来源。
